
\begin{frame}{Outcomes: Three Layers of the Contest}
\hypertarget{src:outcomes}{}
\small
The unit of analysis is the local election. We estimate whether adversarial litigation intensity changes the contest itself; who wins a given seat is not the object. Outcomes fall into three layers, each measured as a 2020--2024 change. Leveling and Barrier predict opposite signs in the first two layers, and the same sign in the third.
\vspace{10pt}
\begin{center}
\renewcommand{\arraystretch}{1.7}
\resizebox{\linewidth}{!}{%
\begin{tabular}{@{}l l c c@{}}
\toprule
 & Outcomes & \textcolor{mygreen}{Leveling predicts} & \textcolor{myred}{Barrier predicts} \\
\midrule
1\ \ Candidate supply & Entrants, women and non-white candidates, party count, field size & broader pool & fewer, safer names \\
2\ \ Voter response & Turnout, blank and null votes & engagement sustained & withdrawal \\
3\ \ Contest outcome & Top-two margin, effective number of candidates, incumbency advantage & ambiguous & consolidation \\
\bottomrule
\end{tabular}}
\end{center}
\vspace{10pt}
{\centering\footnotesize Only layer 2 discriminates. A narrower contest is consistent with either face, so the face is assigned on the ballot.\par}
\end{frame}
